\chapter*{前書き}
\addcontentsline{toc}{chapter}{前書き}

私たちは今、少し不思議な時代を生きています。

スマートフォンの中のAIは、文章を書き、画像を作り、プログラムを組み、難しい質問にもそれらしく答えます。一方で、夜空を見上げれば、宇宙は138億年という途方もない時間をかけて膨張し、星を作り、元素を作り、生命を生み、そしてついにはその宇宙自身について考える知性を生み出したことに気づきます。

AIと宇宙。いっけん遠く離れた二つの主題は、実は深いところで同じ問いに結びついています。

\begin{quote}
複雑なものの背後には、どのようなシンプルな法則が隠れているのか。
\end{quote}

リンゴが落ちる。惑星が回る。熱いコーヒーが冷める。光が波として進む。電子は確率の雲として振る舞う。無数の粒子が集まると温度や圧力が生まれる。巨大な宇宙は膨張し、見えないダークマターやダークエネルギーに支配されている。そして、数千億のパラメータを持つニューラルネットワークは、単純な掛け算と足し算の積み重ねから、言葉や意味を扱うように見える。

本書は、これらを別々の話としてではなく、一つの長い知の系譜として読み解くための本です。

\section*{物理学をAIの言葉で読み直す}

物理学は、自然界の複雑な現象を、できるだけ少ない原理と数式で記述しようとしてきた学問です。ニュートン力学は運動を、マクスウェル方程式は電磁場を、熱力学と統計力学は無数の粒子が生むマクロな秩序を、量子力学はミクロな世界の確率的な振る舞いを説明してきました。

一方、AIもまた、膨大で複雑なデータの中からパターンを見つけ、予測し、圧縮し、表現する技術です。AIの学習は「損失関数」という谷底を探す最適化であり、ニューラルネットワークの内部には高次元の幾何学が広がっています。生成AIはノイズから秩序を作り出し、PINNsは物理法則をニューラルネットワークに埋め込み、AlphaGeometryやAlphaProofは数学的証明の探索に挑みます。

つまり、AIは物理学と無関係な最新技術ではありません。むしろAIは、力学、熱力学、統計力学、量子論、複雑系、情報理論が長い時間をかけて育ててきた考え方を、別の姿で現代に立ち上がらせているのです。

本書では、古典力学の「最適化」から出発し、エントロピー、場、波、相対論、量子論、統計力学、素粒子論、宇宙論へと進みます。そして最後に、AIそのものを新しい自然現象として見つめ直します。これは、物理学の歴史をたどる旅であると同時に、AIを理解するための地図でもあります。

\section*{数式は敵ではなく、世界を圧縮する言葉である}

本書には多くの数式が登場します。けれども、数式を「暗記すべき記号の列」として読む必要はありません。

数式とは、長い説明を短く圧縮した言葉です。例えば、ニュートンの運動方程式は「力を受けた物体は加速度を持つ」という直感を、たった一行で表します。ボルツマンの式は「乱雑さ」と「あり得る状態の数」を結びつけます。シュレディンガー方程式は、ミクロな世界の波がどのように変化するかを記述します。

大切なのは、数式を見た瞬間にすべてを計算できることではありません。その式が、どのような現象を、どのような見方で切り取っているのかをつかむことです。数式は世界を閉じ込める檻ではなく、複雑な世界を見通すための窓です。

本書では、数式の厳密な導出よりも、その背後にある直感と意味を重視します。必要に応じて式を眺め、言葉の説明と行き来しながら読んでください。最初は曖昧に見えても、章を進むにつれて、同じ構造が何度も姿を変えて現れることに気づくはずです。

\section*{この本の読み方}

第1章から第8章までは、近代物理学がどのように生まれ、世界の見方を変えてきたかをたどります。ここでは、力学、最小作用、熱力学、電磁気学、波、相対論、前期量子論が、AIの考え方とどのように響き合うかを見ていきます。

第9章から第13章では、連続体力学、量子力学、統計力学、物質科学、素粒子物理学へと進みます。ここでは、単純な方程式では扱いきれない複雑さ、無数の粒子から生まれる創発、そしてミクロな世界を支配する対称性や場の考え方が中心になります。

第14章から第17章では、視野を宇宙全体とAIそのものへ広げます。宇宙論、数学におけるAI、AIの物理学、そして「理解とは何か」という問いを扱います。ここで本書の主題は、単なる技術解説を超えて、人間の知性とAIの関係へと向かいます。

各章は順番に読むことを想定していますが、興味のある章から入っても構いません。分からない箇所があっても、そこで止まる必要はありません。物理学もAIも、一度で完全に分かるものではなく、何度も同じ景色を違う角度から見ることで、少しずつ輪郭が浮かび上がるものだからです。

\section*{AI時代に、なぜ物理学を学ぶのか}

AIが強力になるほど、「人間は何を理解すべきなのか」という問いは、むしろ鋭くなります。

AIは大量のデータから高精度の予測を出すことができます。しかし、その予測がなぜ成り立つのか、どこまで信頼できるのか、どのような条件で破綻するのかを考えるには、今でも人間の側に深い概念の理解が必要です。物理学は、そのための訓練になります。

物理学を学ぶとは、単に公式を覚えることではありません。複雑な現象の中から本質的な変数を選び、モデルを作り、仮定を置き、限界を見極め、必要ならそのモデルを捨てて新しい見方へ移ることです。これは、AI時代にこそ必要な知的態度です。

AIは答えを出すかもしれません。しかし、どの問いを立てるべきか、その答えにどのような意味を与えるべきか、そしてその理解を人類の知識としてどう位置づけるべきかは、まだ私たち人間に残された仕事です。

本書の目的は、AIを神秘化することでも、物理学を過去の偉人たちの記念碑として眺めることでもありません。AIと物理学を互いの鏡として読み解き、複雑な世界に対して、読者自身がより深く、より自由に考えるための足場を作ることです。

ページをめくるたびに、古典力学の谷底、熱力学のエントロピー、量子力学の確率、宇宙論のダークマター、そしてAIのブラックボックスが、ばらばらの話ではなく、一つの大きな問いの異なる表情として見えてくるはずです。

この本が、その問いへ向かう旅の、確かな出発点になることを願っています。
